\documentclass[exercice,a4]{thedocument}%
\usepackage{texmaths-tikz}%
\usepackage{pstricks}%
\usepackage{pst-bezier}%
\usepackage{pstricks-add}%
\usepackage{pst-eucl}%
\usepackage{pst-func}%
\usepackage{pst-tree}%
\usepackage{pst-circ}%
\usepackage{tikz}%
\usepackage{pgfplots}%
\usepackage{tkz-tab}%
\usepackage{tikzpagenodes}%
\usepackage{tabularray}%
\usepackage{tabularx}%
\usepackage{multirow}%
\usepackage{multicol}%
\usepackage{scratch3}%
\usepackage{enumitem}%
\usepackage[np]{numprint}%
\usepackage{pas-tableur}%
\usepackage{listings}%
\usepackage{graphicx}%
\usepackage{adjustbox}%
\usepackage{needspace}%
\usepackage{xcolor}%
\usepackage{eurosym}%
\usepackage{diagbox}%
\usepackage{colortbl}%
\usepackage{multido}%
\usepackage{icomma}%
\usepackage[french]{babel}%
\startdatakeys
    \source{26GENMATAN1}
    \BO{2026}
    \cycle{4}
    \competencesDuSocle{CA,RA,CO}
    \objectifsApprentissage{C4A1aa,C4A2ad,C4B2cb,C4D2aa,C4D3ac}
    \motsCles{réciproque de Pythagore,aire,droites parallèles,Pythagore,Thalès,angle}
\enddatakeys
\definecolor{gris}{gray}{0.5}%
\providecommand{\wideparen}[1]{\overset{\frown}{#1}}%
\begin{document}%
\providecommand{\degres}{^\circ}%
%-------------------------------------------------------------

La figure ci-contre n'est pas représentée en vraie grandeur.

\vspace{6pt}
\begin{minipage}[t]{0.48\textwidth}
\vspace{0pt}
Les points B, A et E sont alignés.\\
Les points C, A et D sont alignés.\\
Le triangle ABC est rectangle en B.
\begin{itemize}
\item DE $= 4{,}8$ cm
\item AD $= 7{,}3$ cm
\item AE $= 5{,}5$ cm
\item BC $= 7{,}2$ cm.
\end{itemize}
\end{minipage}
\hfill
\begin{minipage}[t]{0.48\textwidth}
\vspace{0pt}
\begin{center}
%début Tikz
%\begin{tikzpicture}[scale=0.9]
%% Coordonnées
%\coordinate (B) at (0,0);
%\coordinate (A) at (1.8,0);
%\coordinate (E) at (4.2,0);
%\coordinate (C) at (1.8,3.5);
%\coordinate (D) at (3.5,-2.0);
%% Triangle ABC
%\draw[thick] (B) -- (A) -- (C) -- cycle;
%% Angle droit en B
%\draw (0.25,0) -- (0.25,0.25) -- (0,0.25);
%% Extension vers E et D
%\draw[thick] (A) -- (E);
%\draw[thick] (A) -- (D);
%% Triangle AED
%\draw[thick] (E) -- (D);
%% Angle droit en E
%\draw ($(E)!0.3cm!(A)$) -- ++(0,-0.3) -- ($(E)+(0,-0.3)$);
%% Labels
%\node[below left] at (B) {B};
%\node[above] at (C) {C};
%\node[above right, xshift=2pt] at (A) {A};
%\node[right] at (E) {E};
%\node[below right] at (D) {D};
%\end{tikzpicture}
%Fin Tikz
%PSTricks
\psset{unit=1cm,arrowsize=2pt 3}
\begin{pspicture}(5.8,7)
\pspolygon(0.4,3)(5.4,3)(5.4,0.2)(0.4,6.2)%BEDC
\uput[ur](3,3){A}\uput[dl](0.4,3){B}\uput[ur](5.4,3){E}\uput[dr](5.4,0.2){D}
\uput[ul](0.4,6.2){C}
\psframe(0.4,3)(0.6,3.2)
\end{pspicture}
%Fin PSTricks
\end{center}
\end{minipage}

\vspace{10pt}

\begin{enumerate}
\item Montrer que le triangle AED est un triangle rectangle en E.
\item Calculer l'aire du triangle AED.
\item Pourquoi peut-on affirmer que les droites (BC) et (ED) sont parallèles ?
\item Calculer la valeur exacte de la longueur AB.
\item On admet que l'angle $\widehat{\text{ACB}}$ mesure environ $49\degres$. En déduire la mesure de l'angle $\widehat{\text{ADE}}$.
\end{enumerate}

\newpage

%-------------------------------------------------------------
\medskip
\showthesource
\end{document}
